\documentclass[12pt]{article}

\usepackage{fullpage}
%\usepackage{amsmath}
%\usepackage{amssymb}
%\usepackage{amsthm}
%\usepackage{amscd}
%\usepackage{setspace}
%\usepackage{wasysym}
%\usepackage{algorithm}
%\usepackage{algpseudocode}

\begin{document}
\title{Review of ``UNIQUE - a Unified framework for wave- based Inversion, Quantification of Uncertainties, and Experimental design", NWO Open Technology Program}
\date{}
\maketitle

This proposal addresses the important challenge of Full Waveform Inversion (FWI), that is, the extraction of valid information on the internal structure of bodies by model-based fitting of remote wave measurements. After decades as a research topic, FWI became commercially important in exploration seismology for petroleum prospecting about 20 years ago, and has become a dominant technology in that field. More recently, a number of academic and commercial research groups have made serious efforts to adapt FWI to biomedical ultrasonic imaging and nondestructive evaluation of manufactured structures.

As has been understood for 40 years, FWI by straightforward least-squares data fitting faces a serious obstacle: if the model-predicted waveforms are seriously out-of-phase with the observed waveforms, then minimization of mean-square misfit can drive the model estimate to a uselessly erroneous apparent local optimum value. This "cycle skipping" phenomenon is not a s well-understood as it should be, but is easily observed in synthetic numerical experiments and occurs in practice as well. A great deal of effort has been devoted to modifying the basic least-squares fit criterion used in early work on FWI, augmenting it with other data less sensitive to cycle-skipping, or replacing it altogether by some substantiallly different measure of error.

The line of work underlying this proposal originates in reference Leeuwen13,

at which point switched to filling our their form.
